Las bases conceptuales fijan el vocabulario del dominio bancario--técnico de este plan de tesis.
Se complementan con las definiciones operativas del Capítulo~1 (§1.1 y §1.4) y
con el mapa semántico unificado de la sección~\ref{sec:glosario}.

\subsection{API REST y contrato OpenAPI}

En este plan, el \textbf{contrato OpenAPI} (YAML/JSON) es el artefacto concreto que el banco
publica para describir su API: \textit{paths}, operaciones, \textit{schemas} y metadatos legibles
por herramientas de documentación y análisis. En banca suele documentar pagos, cuentas o préstamos
\textbf{antes} de exponer el servicio a canales digitales o socios \textit{fintech}.

Casas et al.\ (2021) confirman OpenAPI como estándar de facto en APIs web. La capa REST subyacente
---recursos, verbos HTTP y representaciones serializadas en \textit{schemas}--- se toma como
\textbf{contexto}, no como objeto de estudio: lo que se evalúa es el contrato publicado, no la
arquitectura de despliegue ni el tráfico en producción.

\subsection{Modelo de referencia BIAN}

\textbf{BIAN} organiza capacidades bancarias en \ServiceDomain{}s: unidades de negocio descritas
mediante \Behavior{}s (acciones), \BusinessObject{}s (entidades) y relaciones de control. En este
plan, BIAN actúa como \textbf{modelo de referencia} para contrastar si un contrato OpenAPI bancario
refleja la semántica y la estructura esperadas del dominio (Farzi, 2021). No se evalúan
implementaciones en código ni despliegues productivos.

\subsection{Alineación, coherencia y desalineación}

\textbf{Alineación} designa el grado de correspondencia estructural y semántica entre un contrato
OpenAPI bancario y un \ServiceDomain{} BIAN emparejado. \textbf{Coherencia} se emplea como sinónimo
operativo cuando el \AlignmentScore{} integrado alcanza niveles Altos o Medios según la escala del
modelo. \textbf{Desalineación} agrupa las divergencias detectables; la tipología (estructural,
semántica, de modelado, REST/contrato) se define en el Capítulo~1 (Tabla~\ref{tab:tipologias}).

\subsection{Modelo de puntuación}

El \textbf{modelo de alineación} propuesto descompone la evaluación en tres componentes ---
\SemanticScore{} (S), \StructuralScore{} (E) y \CoverageScore{} (C) --- agregados en
\AlignmentScore{} mediante pesos $\alpha$, $\beta$ y $\gamma$. La \Cref{fig:bases-conceptuales}
resume la relación entre capa técnica (OpenAPI), capa de negocio (BIAN) y producto del modelo.

\begin{figure}[H]
  \centering
  \includegraphics[width=0.92\textwidth]{Figura_05_Bases_Conceptuales}
  \caption{Bases conceptuales: contrato OpenAPI, \ServiceDomain{} BIAN y modelo de alineación S/E/C}
  \label{fig:bases-conceptuales}
\end{figure}
